\begin{figure}[ht]
\centering
\begin{tikzpicture}[x=1cm, y=1cm, font=\small]

% Error versus step size on log-log axes for three derivative schemes
% on a smooth analytic function: forward difference (V apex high),
% central difference (V apex lower), complex-step (monotone fall to a
% machine-precision floor with no round-off arm).
% x-axis: log10(h) from -16 (left) to 0 (right). Map: x = 0.55*(log h +16).
%   log h = -16 -> x = 0; log h = 0 -> x = 8.8.
% y-axis: log10(error) from 0 (top) to -16 (bottom).
%   Map: y = 5.0 + 0.3*log err. log err = 0 -> y = 5.0; -16 -> y = 0.2.

\draw[->, thick] (-0.3, 0) -- (9.4, 0)
  node[right] {$\log_{10} h$};
\draw[->, thick] (0, -0.3) -- (0, 5.6)
  node[above, align=center] {$\log_{10}(\text{error})$};

% x ticks at log h = -16, -12, -8, -4, 0
\foreach \lh in {-16,-12,-8,-4,0} {
  \pgfmathsetmacro\xx{0.55*(\lh + 16)}
  \draw (\xx, 0) -- (\xx, -0.12);
  \node[anchor=north, font=\scriptsize] at (\xx, -0.12) {\lh};
}
% y ticks at log err = 0, -4, -8, -12, -16
\foreach \le in {0,-4,-8,-12,-16} {
  \pgfmathsetmacro\yy{5.0 + 0.3*\le}
  \draw (0, \yy) -- (-0.12, \yy);
  \node[anchor=east, font=\scriptsize] at (-0.14, \yy) {\le};
}

% Forward difference: trunc ~ h (slope +1 on right), round-off ~ eps/h
% (slope -1 on left), apex near log h = -8, error ~ -8.
% Build from sample points (log h, log err).
\draw[very thick, engred, dotted]
  plot coordinates {
    ({0.55*(-16+16)}, {5.0 + 0.3*(0)})
    ({0.55*(-12+16)}, {5.0 + 0.3*(-4)})
    ({0.55*(-8+16)},  {5.0 + 0.3*(-8)})
    ({0.55*(-6+16)},  {5.0 + 0.3*(-6)})
    ({0.55*(-2+16)},  {5.0 + 0.3*(-2)})
    ({0.55*(0+16)},   {5.0 + 0.3*(0)})
  };
\node[engred, anchor=west, font=\scriptsize] at (5.2, 4.7)
  {forward diff.};

% Central difference: trunc ~ h^2 (slope +2 on right), round-off ~ eps/h
% (slope -1), apex near log h = -5.3, error ~ -11.
\draw[very thick, engamber, dashed]
  plot coordinates {
    ({0.55*(-16+16)}, {5.0 + 0.3*(-1)})
    ({0.55*(-12+16)}, {5.0 + 0.3*(-5)})
    ({0.55*(-8+16)},  {5.0 + 0.3*(-9)})
    ({0.55*(-5.3+16)},{5.0 + 0.3*(-11)})
    ({0.55*(-2+16)},  {5.0 + 0.3*(-4.6)})
    ({0.55*(0+16)},   {5.0 + 0.3*(-0.6)})
  };
\node[engamber, anchor=west, font=\scriptsize] at (5.2, 4.3)
  {central diff.};

% Complex step: trunc ~ h^2 (slope +2 on right), no round-off arm; the
% error falls until it hits the machine floor ~ -16 and stays there.
\draw[very thick, engblue]
  plot coordinates {
    ({0.55*(-16+16)}, {5.0 + 0.3*(-16)})
    ({0.55*(-12+16)}, {5.0 + 0.3*(-16)})
    ({0.55*(-8+16)},  {5.0 + 0.3*(-16)})
    ({0.55*(-6+16)},  {5.0 + 0.3*(-12.6)})
    ({0.55*(-2+16)},  {5.0 + 0.3*(-4.6)})
    ({0.55*(0+16)},   {5.0 + 0.3*(-0.6)})
  };
\node[engblue, anchor=west, font=\scriptsize] at (5.2, 0.55)
  {complex step};

% machine floor guide line
\draw[thin, enggray, densely dotted] (0, {5.0 + 0.3*(-16)})
  -- (8.8, {5.0 + 0.3*(-16)});

\end{tikzpicture}
\caption[Finite-difference and complex-step accuracy against step
size]{Absolute error of three first-derivative schemes on a smooth
analytic function, as the step size $h$ sweeps from $10^{0}$ down to
$10^{-16}$. The forward and central differences each show the
characteristic V: a truncation arm that falls as a power of $h$ on the
right and a round-off arm that rises as $\varepsilon/h$ on the left,
with the central scheme's apex one order lower because its truncation
falls as $h^{2}$ rather than $h$. The complex-step derivative
$\operatorname{Im} f(x + ih)/h$ has no subtraction of nearly equal
numbers, so it carries no round-off arm: the error falls as $h^{2}$
and then sits on the machine-precision floor for every smaller $h$.
The step-size tuning problem that dominates finite differencing
disappears for the complex-step method wherever the function can be
evaluated in complex arithmetic.}
\label{fig:vol02:ch05:complex-step}
\end{figure}
